\chapter{Mock Interviews}\label{ch:iv:mock-interviews}

Two transcripts of the same forty-minute interview for the same trading role, with the same questions and the same interviewer.
Both candidates reached most of the right answers. One was hired. The assessor's margin notes on the other say ``did not update'',
``did not check'', and, twice, ``no number''. This chapter prints six complete interviews, one for each column of the table of
\cref{ch:iv:what-each-interview-tests-by-role}, as a candidate answered them, with an assessor's commentary in the margin of
each exchange. The candidates are composites written for the book; the questions are new, and the solutions give the answer the
assessor was looking for.

\section{How to read a transcript}

\begin{definition}[Mock interview]\label{def:iv:mock-interviews:mock}
A \emph{mock interview} is a practice interview run under the conditions of a real one (the same format, time and kind of
questions, an interviewer who scores against a rubric) and followed by feedback on each answer.
\end{definition}

Each transcript gives the setting, then the exchanges: \textsc{Interviewer} and \textsc{Candidate} lines, with the assessor's
notes in italics. The questions the interviewer asks are numbered like every question in the book, and their model answers are in
the solutions. Read a transcript twice: first as the candidate, stopping at each question to answer it yourself; then as the
assessor, asking what each answer lets the assessor write on the rubric. The assessor notes the same few things in every
interview: a number given, a check made, an update on new information, an explanation of the method, and composure after a mistake.
Each question is followed by the interviewer's follow-up, because that is where most of the scoring happens: the first answer shows
what the candidate prepared, the follow-up what the candidate can do. The trader's interview is printed twice, as the hook promised,
once for the candidate who was hired and once for the one who was not; the last section collects the seven scorecards.

\section{Trader at a market maker}

\emph{Setting.} A final-round interview for a graduate trader at an options market maker; forty minutes; an experienced trader
interviews; mental arithmetic, a market and a betting question.

\begin{interviewq}[$\star$ \iqroles{trader} \iqfirm{market maker}]\label{iq:iv:mock-interviews:1}
``Quickly: fifteen per cent of 64\,000.''
\end{interviewq}

\noindent\textsc{Candidate.} ``Ten per cent is 6\,400, five per cent is 3\,200, so 9\,600.''

\noindent{\small\itshape\color{omIq}Assessor: three seconds, the method said in one breath. Top anchor on speed.}

\noindent\textsc{Interviewer.} ``Take seven per cent off that.''

\noindent\textsc{Candidate.} ``Seven per cent of 9\,600 is 672, so 8\,928. Check: 0.93 of 9\,000 is 8\,370 and 0.93 of 600 is 558;
together 8\,928.''

\noindent{\small\itshape\color{omIq}Assessor: second route unprompted.}

\begin{interviewq}[$\star\star$ \iqroles{trader} \iqfirm{market maker}]\label{iq:iv:mock-interviews:2}
``I am going to toss a coin twenty times. I have already tossed it five times and seen three heads, and you have not. Make me a
market on the total number of heads.''
\end{interviewq}

\noindent\textsc{Candidate.} ``Fifteen tosses to go at a half each, seven and a half, plus your three: ten and a half. Standard
deviation about two. I'll quote 10 at 11.''

\noindent\textsc{Interviewer.} ``Buy at 11. Again.''

\noindent\textsc{Candidate.} ``You know three heads have happened, and so do I, so your buying tells me nothing about the tosses to
come. I'll stay at 10 at 11, and I'm now short one at 11 against a fair value of 10.5, which I like.''

\noindent{\small\itshape\color{omIq}Assessor: centre right, width tied to the spread, and the key point: the counterparty's
information was already public, so no update. Said the position and its value unprompted.}

\noindent\textsc{Interviewer.} ``I toss the sixth. Heads. Your market?''

\noindent\textsc{Candidate.} ``Four heads in six and fourteen to go: $4 + 7 = 11$, standard deviation $\sqrt{3.5} \approx 1.87$. I'm
10.5 at 11.5. And my short at 11 is now worth nothing in expectation: that toss cost me half a point.''

\noindent{\small\itshape\color{omIq}Assessor: this time the information was real, and the market moved at once; marked the position to the
new fair value without being asked. Top anchor on updating.}

\begin{interviewq}[$\star\star\star$ \iqroles{trader, risk} \iqfirm{market maker}]\label{iq:iv:mock-interviews:3}
``A deck has ten red and six black cards left. You may bet any part of a bankroll of 400 at even money that the next card is red,
once. How much do you bet if you want to maximise long-run growth over many such games, and why not everything?''
\end{interviewq}

\noindent\textsc{Candidate.} ``The chance is ten sixteenths, 62.5\%. At even money Kelly says two $p$ minus one, a quarter, so 100.
Not everything, because six times in sixteen I lose it all, and the log of zero is minus infinity\ldots\ sorry, I should say: betting
everything maximises the expected bankroll but not its growth.''

\noindent{\small\itshape\color{omIq}Assessor: correct fraction and stake; corrected his own phrasing, which I note as a positive.
Could have given the growth rate per game (about 3\% of log-wealth) and mentioned fractional Kelly if the 62.5\% were uncertain.}

\noindent\textsc{Interviewer.} ``Suppose you are not sure of the count: the chance of red could be anywhere from 55\% to 70\%.''

\noindent\textsc{Candidate.} ``Kelly at 55\% is a tenth of the bankroll, 40; at 70\% it is four tenths, 160. The mistakes aren't
symmetric: if the truth is 55\% and I bet a quarter, my growth rate is negative, about $-0.7\%$ a game; if the truth is 70\% and I bet a
tenth, I still grow, at 3.5\% against the best possible 8.2\%. So I lean low: 40 to 50.''

\noindent{\small\itshape\color{omIq}Assessor: computed both ends and the asymmetry, and turned it into a number. Strong hire.}

\subsection*{The same interview, the candidate who was not hired}

\noindent\textsc{Interviewer.} ``Quickly: fifteen per cent of 64\,000.''

\noindent\textsc{Candidate.} ``About ten thousand. \ldots\ 9\,600.''

\noindent{\small\itshape\color{omIq}Assessor: guessed, then computed; fine, slower.}

\noindent\textsc{Interviewer.} (The coin, as above.) ``Make me a market on the total.''

\noindent\textsc{Candidate.} ``Around ten. 8 at 12.''

\noindent{\small\itshape\color{omIq}Assessor: no number for the uncertainty; a width of four is two standard deviations, which is wide for
a coin.}

\noindent\textsc{Interviewer.} ``Buy at 12. I toss the sixth: heads.''

\noindent\textsc{Candidate.} ``I'll stay at 8 at 12.''

\noindent{\small\itshape\color{omIq}Assessor: did not update: the toss moved fair value from 10.5 to 11, and the candidate is now short one
at 12 without saying so.}

\noindent\textsc{Interviewer.} (The deck, as above.) ``How much do you bet?''

\noindent\textsc{Candidate.} ``A good part of it, since red is more likely. Most of it, but not everything, to be safe.''

\noindent{\small\itshape\color{omIq}Assessor: no number.}

\noindent\textsc{Interviewer.} ``How much, exactly?''

\noindent\textsc{Candidate.} ``Three hundred. It feels right: red is much more likely than black.''

\noindent{\small\itshape\color{omIq}Assessor: did not check: at 300 a loss leaves 100, and the growth rate is $0.625\ln 1.75 + 0.375\ln 0.25
\approx -0.17$ a game, so the bankroll shrinks over many games. Never asked what ``long-run growth'' meant.}

Both candidates reached most of the right answers. What separates them is visible in the notes: numbers given without being asked,
checks made aloud, and positions updated when, and only when, the information changed.

\section{Researcher at a systematic fund}

\emph{Setting.} A \omterm{def:iv:what-each-interview-tests-by-role:loop}{technical interview} for a quantitative researcher at a systematic equity fund; forty-five minutes; a senior
researcher interviews; statistics, a regression puzzle, and the design of a test.

\begin{interviewq}[$\star$ \iqroles{researcher} \iqfirm{systematic fund}]\label{iq:iv:mock-interviews:4}
``A daily signal has an information coefficient of 0.02 over ten years. Is that significant?''
\end{interviewq}

\noindent\textsc{Candidate.} ``Ten years is about 2\,500 days. As a correlation on 2\,500 independent observations, the $t$-statistic
is about 0.02 times the square root of 2\,500, so 1. Not significant. But a cross-sectional coefficient of 0.02 can still be
valuable across many stocks; I'd want to know the breadth and the turnover before dismissing it.''

\noindent{\small\itshape\color{omIq}Assessor: fast and right, and the ``but'' shows she knows the fundamental law. Did not ask whether
the observations are autocorrelated, which would make the $t$ smaller still.}

\noindent\textsc{Interviewer.} ``The daily coefficients have a first-order autocorrelation of 0.5. Now?''

\noindent\textsc{Candidate.} ``For an autoregression of coefficient $\rho$ the variance of the mean is inflated by about $(1 + \rho)/(1 -
\rho)$, here 3, so the $t$-statistic falls by $\sqrt3$, to about 0.58.''

\noindent{\small\itshape\color{omIq}Assessor: had the formula ready.}

\begin{interviewq}[$\star\star$ \iqroles{researcher} \iqfirm{systematic fund}]\label{iq:iv:mock-interviews:5}
``You regress next-month returns on a value signal and get a positive coefficient. Adding sector dummies makes it negative. Which do
you believe?''
\end{interviewq}

\noindent\textsc{Candidate.} ``Neither yet. Without sectors, the value coefficient partly measures which sectors are cheap, and
those sectors did well. Within sectors, cheap stocks did worse. The two answer different questions: a sector-neutral strategy should
believe the second one. I'd check whether the within-sector result is stable over time and whether one sector drives it.''

\noindent{\small\itshape\color{omIq}Assessor: named the omitted-variable mechanism without being prompted and connected each
coefficient to a strategy. Top anchor.}

\noindent\textsc{Interviewer.} ``One sector drives the within-sector result.''

\noindent\textsc{Candidate.} ``Then it's a statement about that sector, not about value. I'd report it as such, re-estimate without it,
and ask whether there is a reason, an accounting convention for example, that makes value measure something different there.''

\noindent{\small\itshape\color{omIq}Assessor: good instinct to look for a mechanism before a trade.}

\begin{interviewq}[$\star\star\star$ \iqroles{researcher} \iqfirm{systematic fund}]\label{iq:iv:mock-interviews:6}
``Design a test of whether a stock's overnight return predicts its next day's intraday return, on 500 stocks over ten years.''
\end{interviewq}

\noindent\textsc{Candidate.} ``That's 1.26 million stock-days, but only about 2\,500 days, and returns on the same day are
correlated, so I'd cluster standard errors by day or run a daily cross-sectional regression and test the mean of the daily
coefficients. Overnight return from the close to the open, intraday from the open to the close, measured at prices one could
trade, so I'd use the opening auction price. Then costs, because the signal trades every stock every day. And I'd tell you in
advance what result would make me drop it.''

\noindent\textsc{Interviewer.} ``And if you had tried five definitions of `overnight' before this one?''

\noindent\textsc{Candidate.} ``Then I report five and correct for them; the best of five is biased up.''

\noindent\textsc{Interviewer.} ``The best of the five had $t = 2.2$.''

\noindent\textsc{Candidate.} ``Two-sided, $p \approx 0.028$; with a Bonferroni correction for five, about 0.14. Not significant. I'd
keep the one definition I would have chosen first, and test it on the next year.''

\noindent{\small\itshape\color{omIq}Assessor: the effective sample size, the clustering, the tradable prices, the costs, and the
pre-registered kill criterion. Strong hire on this interview.}

\section{Developer at a proprietary firm}

\emph{Setting.} A \omterm{def:iv:what-each-interview-tests-by-role:loop}{technical interview} for a software engineer on a trading-systems team; forty-five minutes; a senior engineer
interviews; complexity, a coding question and a small design.

\begin{interviewq}[$\star$ \iqroles{developer} \iqfirm{proprietary firm}]\label{iq:iv:mock-interviews:7}
``How do you find whether any client order identifier repeats in a day's log of ten million orders, and what does it cost?''
\end{interviewq}

\noindent\textsc{Candidate.} ``A hash set: insert each identifier and stop at the first one already there. Linear time, and memory
for up to ten million identifiers, hundreds of megabytes in a naive set. If memory were tight I'd sort and scan, $n \log n$,
or use a Bloom filter first and check its hits.''

\noindent{\small\itshape\color{omIq}Assessor: correct, with the memory cost and two alternatives. Good.}

\noindent\textsc{Interviewer.} ``Size the Bloom filter for a 1\% false-positive rate.''

\noindent\textsc{Candidate.} ``$m = -n\ln p/(\ln 2)^2$: $10^7 \times 4.6/0.48 \approx 96$ million bits, 12 megabytes, with $(m/n)\ln 2
\approx 6.6$, so seven hash functions. About 100\,000 false alarms to check exactly, which is cheap.''

\noindent{\small\itshape\color{omIq}Assessor: had the formula, and turned the rate into a count.}

\begin{interviewq}[$\star\star$ \iqroles{developer} \iqfirm{proprietary firm}]\label{iq:iv:mock-interviews:8}
``Write a cache of instrument definitions of fixed capacity that evicts the least recently used entry, with constant-time get and
put.''
\end{interviewq}

\noindent\textsc{Candidate.} ``A hash map from key to a node in a doubly linked list ordered by recency; get moves the node to the
front; put inserts at the front and, if over capacity, removes the tail. In Python an ordered dictionary does both jobs.'' (Writes
it; tests it on a capacity of two: put A, put B, get A, put C, and checks that B was evicted.)

\noindent{\small\itshape\color{omIq}Assessor: right structure, wrote a test before being asked. Top anchor on testing.}

\noindent\textsc{Interviewer.} ``Two threads share it.''

\noindent\textsc{Candidate.} ``A get also writes, since it moves the key, so a reader-writer lock buys nothing. One mutex around both
operations first; if a profile shows contention, shard the cache by key hash with a mutex per shard, which weakens `least recently
used' to per-shard, and I'd say so in the interface.''

\noindent{\small\itshape\color{omIq}Assessor: spotted that reads mutate; named the semantic cost of sharding.}

\begin{interviewq}[$\star\star\star$ \iqroles{developer} \iqfirm{proprietary firm}]\label{iq:iv:mock-interviews:9}
``The venue allows 100 messages a second with bursts of up to 20. Design the throttle in our gateway. If 50 messages arrive at once
and then one every ten milliseconds for a second, starting half a second later, how many of the 150 go out without waiting?''
\end{interviewq}

\noindent\textsc{Candidate.} ``A token bucket: capacity 20, refilled at 100 a second; a message goes out if a token is available,
otherwise it waits or is rejected, which is a policy decision I'd ask the traders about. At time zero 20 go and 30 wait. After half
a second the bucket is full again at 20, and new messages arrive at exactly the refill rate, so all 100 go. That's 120 without
waiting, if the 30 are rejected.''

\noindent\textsc{Interviewer.} ``And if they queue?''

\noindent\textsc{Candidate.} ``They leave one every ten milliseconds, the last at 0.3 seconds, and the bucket refills to 20 by 0.5, so
the second wave is untouched: all 150 go, 30 of them late, the last by 300 milliseconds. Whether that's better is the desk's call: an
order 300 milliseconds old can be worse than no order. I'd reject new orders and tell the strategy, and queue cancels, or give them
their own budget if the venue allows.''

\noindent{\small\itshape\color{omIq}Assessor: correct model and counts in both policies, and turned the policy into a business question.
Hire.}

\section{Machine-learning engineer at a systematic fund}

\emph{Setting.} A \omterm{def:iv:what-each-interview-tests-by-role:loop}{technical interview} for a machine-learning engineer who will put research models into production; forty
minutes; an engineering lead interviews.

\begin{interviewq}[$\star$ \iqroles{mle} \iqfirm{systematic fund}]\label{iq:iv:mock-interviews:10}
``How do you split five years of daily data into training, validation and test sets?''
\end{interviewq}

\noindent\textsc{Candidate.} ``By time: the first three years to train, the fourth to validate, the fifth to test, and a gap
between them as long as the label horizon so that labels do not overlap across the boundary. Or walk forward with expanding
windows. Never shuffle.''

\noindent{\small\itshape\color{omIq}Assessor: the embargo gap mentioned unprompted. Good.}

\noindent\textsc{Interviewer.} ``The labels are five-day forward returns, and you want several test years, not one.''

\noindent\textsc{Candidate.} ``Then a five-day gap at every boundary, and walk forward: at least two years of training, one-year test
windows, so three folds in five years, each trained on everything before it. I'd report the three test years separately as well as
pooled; if one year carries the result, that's the finding.''

\noindent{\small\itshape\color{omIq}Assessor: sized the gap from the label and counted the folds.}

\begin{interviewq}[$\star\star$ \iqroles{mle} \iqfirm{systematic fund}]\label{iq:iv:mock-interviews:11}
``A classifier flags one per cent of events as anomalies, and its accuracy is 99\%. Is it good?''
\end{interviewq}

\noindent\textsc{Candidate.} ``Predicting `normal' always gives 99\% too, so accuracy says nothing. I'd look at precision and
recall on the anomalies, or the precision-recall curve, and at the cost of each kind of error, which decides the threshold.''

\noindent{\small\itshape\color{omIq}Assessor: correct and quick; tied the threshold to costs.}

\noindent\textsc{Interviewer.} ``Yours reaches 80\% recall at a false-positive rate of 5\%. What's its precision?''

\noindent\textsc{Candidate.} ``Per 10\,000 events: 100 anomalies, 80 caught; 9\,900 normal, 495 flagged. Precision is 80 out of 575,
about 14\%: six false alarms for each real one. Whether that's usable depends on who reads the alerts.''

\noindent{\small\itshape\color{omIq}Assessor: counted instead of reaching for a formula; the base rate carried through.}

\begin{interviewq}[$\star\star\star$ \iqroles{mle, developer} \iqfirm{systematic fund}]\label{iq:iv:mock-interviews:12}
``A model's input feature was spread evenly over four bins in training; this week the shares are 10, 20, 30 and 40\%. Compute the
population stability index and tell me what you would do.''
\end{interviewq}

\noindent\textsc{Candidate.} ``PSI is the sum of actual minus expected times the log of their ratio: $-0.15 \ln 0.4$, then $-0.05
\ln 0.8$, then $0.05 \ln 1.2$, then $0.15 \ln 1.6$\ldots\ about 0.23. The usual rule of thumb treats above 0.2 as a large shift. I'd
check whether the change is real (a new regime) or a pipeline fault (a unit change, a stale source), and if real, look at the
model's recent performance before retraining.''

\noindent{\small\itshape\color{omIq}Assessor: computed it, gave the number, and separated a data fault from a real shift before
reaching for retraining (One Quant Book~12, chapter~27). Hire.}

\noindent\textsc{Interviewer.} ``It turns out the vendor switched that feature from per cent to basis points.''

\noindent\textsc{Candidate.} ``Then the model is fine and the pipeline isn't. Fix the unit, rescore the affected days, and add a check on
the feature's range and distribution at ingestion so that the next change pages someone instead of feeding the model garbage.''

\noindent{\small\itshape\color{omIq}Assessor: fixed the class of fault, not the instance.}

\section{Bank quant at a bank}

\emph{Setting.} A \omterm{def:iv:what-each-interview-tests-by-role:loop}{technical interview} for a desk quant on an equity-derivatives desk; forty-five minutes; a senior quant interviews;
pricing and sensitivities.

\begin{interviewq}[$\star$ \iqroles{bank} \iqfirm{bank}]\label{iq:iv:mock-interviews:13}
``A stock at 100 pays a continuous dividend yield of 1\%; the rate is 3\%. What is the two-year forward?''
\end{interviewq}

\noindent\textsc{Candidate.} ``$100 e^{(0.03 - 0.01) \times 2} = 100 e^{0.04}$, about 104.08.''

\noindent{\small\itshape\color{omIq}Assessor: instant. Fine.}

\noindent\textsc{Interviewer.} ``Instead of the yield, a dividend of 2 paid in six months.''

\noindent\textsc{Candidate.} ``Take the dividend's present value off spot, $2e^{-0.015} \approx 1.97$, and grow the rest: $98.03 \times
e^{0.06} \approx 104.09$. Close to the yield case, as it should be: 1\% a year on 100 over two years is about 2.''

\noindent{\small\itshape\color{omIq}Assessor: right treatment and a consistency check.}

\begin{interviewq}[$\star\star$ \iqroles{bank, trader} \iqfirm{bank}]\label{iq:iv:mock-interviews:14}
``A call struck at the money, one year to expiry, stock at 100, 20\% volatility, zero rates: what is its vega per volatility
point?''
\end{interviewq}

\noindent\textsc{Candidate.} ``Vega is $S\varphi(d_1)\sqrt T$. $d_1$ is 0.1, $\varphi(0.1)$ about 0.397, so about 39.7 per unit of
volatility, 0.40 per point. That's also about the at-the-money call price divided by the volatility, since the price is roughly
linear in volatility at the money: 8 over 20.''

\noindent{\small\itshape\color{omIq}Assessor: formula, number and an independent check. Top anchor.}

\noindent\textsc{Interviewer.} ``And for the three-month option?''

\noindent\textsc{Candidate.} ``At the money, vega scales like the square root of time: half, about 0.20 a point.''

\noindent{\small\itshape\color{omIq}Assessor: scaling law rather than recomputation.}

\begin{interviewq}[$\star\star\star$ \iqroles{bank} \iqfirm{bank}]\label{iq:iv:mock-interviews:15}
``Price a one-year digital call struck at 100, same stock, by a call spread, and tell me how big the replication error is and what
the desk should charge for it.''
\end{interviewq}

\noindent\textsc{Candidate.} ``Buy the 99 call and sell the 101 call, scaled by a half: in the Black--Scholes model that gives
about 0.460, the same as $\Phi(d_2)$ to four decimals, because the error falls like the square of the half-width. Only if the
desk can trade exactly those strikes, and the payoff at expiry is not a digital: between 99 and 101 it pays a fraction. The desk
over-hedges, buying the spread on the side that makes the payoff at least the digital's, and charges that difference, which is where
the skew matters.''

\noindent{\small\itshape\color{omIq}Assessor: correct price, the convergence order, and the over-hedge as the practical answer (One
Quant Book~5, chapter~15). Did not quantify the skew adjustment; acceptable at this level.}

\noindent\textsc{Interviewer.} ``The client wants a million on the digital. How many spreads?''

\noindent\textsc{Candidate.} ``Each 99--101 spread pays at most 2, so a million over 2: 500\,000 spreads, bought on the side that
over-hedges, and the desk must be able to trade those strikes in that size.''

\noindent{\small\itshape\color{omIq}Assessor: sized it and added the liquidity caveat.}

\section{Portfolio-manager hire at a multi-manager fund}

\emph{Setting.} A meeting with a candidate portfolio manager at a multi-manager platform; sixty minutes; the head of a business
unit and a risk manager interview; the candidate has brought two years of monthly returns.

\begin{interviewq}[$\star$ \iqroles{trader, researcher} \iqfirm{multi-manager fund}]\label{iq:iv:mock-interviews:16}
``From your monthly returns, what are your annualised Sharpe ratio and your maximum drawdown?''
\end{interviewq}

\noindent\textsc{Candidate.} ``About 1.05 and about 2.4\% from peak, on twenty-four months, after costs, on the capital I ran. Two
years is short: the standard error on that Sharpe ratio is around 0.9.''

\noindent{\small\itshape\color{omIq}Risk manager: precise, and volunteered the uncertainty. We will recompute from the administrator's
statements.}

\noindent\textsc{Risk manager.} ``Take out your best month.''

\noindent\textsc{Candidate.} ``My best month was 2.1\%. Without it the Sharpe ratio is about 0.86. The drawdown doesn't change: the best month
came just before it, so taking it out lowers the peak and the trough together.''

\noindent{\small\itshape\color{omIq}Risk manager: knew the answer, which means the candidate had asked the question first.}

\begin{interviewq}[$\star\star$ \iqroles{trader, researcher} \iqfirm{multi-manager fund}]\label{iq:iv:mock-interviews:17}
``Your gross Sharpe ratio is 1.2 at 8\% volatility with a turnover of twenty times a year and costs of 5 basis points a turn. What
happens to the net Sharpe ratio if we give you ten times the capital?''
\end{interviewq}

\noindent\textsc{Candidate.} ``Today costs take $20 \times 5$ basis points, 1\%, from a gross return of 9.6\%: net 8.6\%, Sharpe
about 1.08. If impact grows with the square root of size, costs rise by about three times at ten times the capital, to about 3.2\%,
so the net Sharpe falls to about 0.8. I'd rather start at three times and measure.''

\noindent{\small\itshape\color{omIq}Business head: understands capacity and proposed the staged increase himself.}

\noindent\textsc{Business head.} ``At what size does your net Sharpe ratio reach 0.5?''

\noindent\textsc{Candidate.} ``At 8\% volatility the net return must be 4\%: $9.6 - \sqrt{m} = 4$ with costs of 1\% today, so $\sqrt m =
5.6$ and $m \approx 31$ times today's capital, if the square-root law held that far, which I doubt. I'd treat ten times as the
ceiling until I had measured it.''

\noindent{\small\itshape\color{omIq}Business head: solved it, then refused to believe the extrapolation. Good.}

\begin{interviewq}[$\star\star\star$ \iqroles{trader, risk} \iqfirm{multi-manager fund}]\label{iq:iv:mock-interviews:18}
``Our rule stops a book out at a 7.5\% drawdown. If your returns were a drifting Brownian motion with a 10\% expected return and 10\%
volatility a year, what is the chance you are stopped out within your first year?''
\end{interviewq}

\noindent\textsc{Candidate.} ``The first-passage probability for Brownian motion with drift: $\Phi$ of minus the barrier minus the
drift, over the volatility, plus $e^{-2\mu a/\sigma^2}$ times $\Phi$ of minus the barrier plus the drift, over the volatility.
With these numbers roughly 17\%. That is measured from the start. If your rule measures from the running peak, and most do, it
is much higher; I'd simulate it, and I'd guess near a half. So between one in six and one in two of people like me would be
stopped by luck in year one; that is worth discussing before I sign.''

\noindent{\small\itshape\color{omIq}Risk manager: correct formula and number, and the distinction between drawdown from the start
and from the peak. The last sentence is the kind of negotiation we respect (One Quant Book~16, chapter~8).}

\noindent\textsc{Risk manager.} ``And if your Sharpe ratio were 2, same volatility?''

\noindent\textsc{Candidate.} ``Drift 20\%: $\Phi(-2.75) + e^{-3}\Phi(1.25) \approx 0.003 + 0.045$, about 5\% from the start. The stop
bites much less often on a better strategy, which is the point of it.''

\noindent{\small\itshape\color{omIq}Risk manager: recomputed in thirty seconds.}

\section{The seven scorecards}

Each interviewer scored the rubric's five rows from 1 to 4: numbers (a value given when one was needed), checks (a second route, a
limit, a units test), updating (moving on information, and only on information), explanation (the method said so another person
could follow it), and composure (after an error or a challenge). The table collects the scores.

\begin{center}
\footnotesize
\setlength{\tabcolsep}{4pt}
\begin{tabular}{lcccccl}
\toprule
Interview & Numbers & Checks & Updating & Explanation & Composure & Decision \\
\midrule
Trader, hired & 4 & 3 & 4 & 4 & 4 & hire \\
Trader, not hired & 2 & 1 & 1 & 2 & 3 & no hire \\
Researcher & 4 & 3 & 4 & 4 & 4 & hire \\
Developer & 4 & 4 & 3 & 4 & 4 & hire \\
ML engineer & 4 & 4 & 4 & 3 & 4 & hire \\
Bank quant & 4 & 4 & 3 & 4 & 4 & hire \\
Portfolio-manager hire & 4 & 3 & 4 & 4 & 4 & hire, if verified \\
\bottomrule
\end{tabular}
\end{center}

Three things are common to the strong transcripts. The first answer is short and has a number in it. The follow-up, which is
where the interviewer learns most, is answered with a new computation rather than a restatement. And every candidate who was
hired said at least once what would change their answer: the sixth toss, the uncertain deck, the five definitions, the queueing
policy, the vendor's units, the skew, the peak. The candidate who was not hired gave most of the same answers eventually; what
the rubric recorded was that the numbers came late, the checks never, and the update not at all.

\begin{omsources}
\item The six candidates are composites written for this book; the questions are new.
\item One Quant Book~5, chapter~15; One Quant Book~7, chapter~28; One Quant Book~12, chapter~27; One Quant Book~16, chapters 3
and~8; One Quant Book~17, chapters 16--22.
\end{omsources}
